\subsection{The obstruction to the rank-zero CM route at exponent seventeen}

Let $F_{17}$ be the smooth projective Fermat curve
$X^{17}+Y^{17}=Z^{17}$, rationally isomorphic to $X_{17}$ by
changing the sign of $Z$, and put
\[
 C_s:\quad v^{17}=u(1-u)^s,\qquad
 A_s=\operatorname{Jac}(C_s),\qquad 1\leq s\leq15.
\]
The genus of each $C_s$ is eight.  The hyperelliptic quotient $C_1$
is isomorphic over $\mathbb Q$ to $y^2+y=x^{17}$.

\begin{proposition}[The degree-eight Hilbert route has zero central values]
\label{prop:lm-seventeen-obstruction}
There is a $\mathbb Q$-isogeny
\[
 \operatorname{Jac}(F_{17})\sim_{\mathbb Q}
 A_1^3\times A_2^6\times A_3^6.
\]
Each $A_i$ is absolutely simple and has CM by
$K=\mathbb Q(\zeta_{17})$.  Its associated weight-two Hilbert
representation over $F=K^+$ has trivial finite central character and
is supercuspidal at the unique place of $F$ above $17$; its conductor
at that place has exponent three.  Thus the nonprincipal-series
alternative in Nekov\'a\v r's 2012 Theorem~B applies to these forms.
Nevertheless,
\[
 L(A_1/\mathbb Q,1)=L(A_2/\mathbb Q,1)
                  =L(A_3/\mathbb Q,1)=0.
\]
In particular, none of these factors supplies the nonzero central
value required by that rank-zero theorem.
\end{proposition}

\begin{proof}
The quotient maps $(x,y)\mapsto(x^{17},xy^s)$ give the rational
isogeny to $\prod_{s=1}^{15}A_s$; see \citet[\S2, pp.~207--208]{GrossRohrlich1978}.  Permuting the triple $(1,s,16-s)$ and
rescaling modulo $17$ gives exactly the following orbits:
\[
 \begin{array}{c|c|c}
 \text{representative}&s\text{-values}&\text{multiplicity}\\ \hline
 (1,1,15)&1,8,15&3\\
 (1,2,14)&2,5,7,9,11,14&6\\
 (1,3,13)&3,4,6,10,12,13&6.
 \end{array}
\]
These equivalences are rational birational equivalences of the
quotient curves.  The displayed product follows.
The CM simplicity assertion follows from \citet[Lemma~1.6 and the following remark, pp.~206--207]{GrossRohrlich1978}: their only
imprimitive case requires a nonidentity cube root of unity modulo
the exponent, and $3\nmid16$.  Equivalently, the cyclic field
$K$ has no proper CM subfield: every nonidentity subgroup of its
cyclic group of order sixteen contains complex conjugation.

Here is an explicit calculation of the local condition.  For
$(r,s,t)=(1,i,16-i)$, set
\[
 z=r^r s^s(t-17)^t\in(\mathbb Z/17^2\mathbb Z)^\times,
 \qquad d=\frac{z^{16}-1}{17}\pmod {17}.
\]
The three values of $(z,d)$, in the order above, are
\[
 (178,8),\qquad(76,13),\qquad(215,11).
\]
In particular, all three $d$ are nonzero.  \citet[Proposition~3.12(2), author p.~16]{ShuExpanded2021} therefore gives conductor
$\mathfrak P^2$ for the Jacobi-sum Hecke character at the prime
$\mathfrak P$ of $K$ over $17$.  Its restriction to
$U^1/U^2\simeq\mathbb F_{17}$ is nontrivial.  With
$\lambda=\zeta_{17}-\zeta_{17}^{-1}$, complex conjugation takes
$\lambda$ to $-\lambda$, and hence acts by $-1$ on $U^1/U^2$.
The character and its conjugate thus have inverse, distinct,
order-seventeen restrictions.  Their induction from $K_{\mathfrak P}$
to $F_{\mathfrak p}$ is an irreducible two-dimensional Weil
representation, giving a supercuspidal local representation.
The tame quadratic extension has different exponent one and residue
degree one; the induction formula for the Artin conductor gives
$1+2=3$.

The central character is also explicit.  In unitary normalization
the Jacobi-sum character restricts to the quadratic character
$\eta_{K/F}$ on $\mathbb A_F^\times$; see \citet[\S3.4, author p.~15]{ShuExpanded2021}.  The determinant character of quadratic automorphic
induction is this restriction times $\eta_{K/F}$, hence is trivial.
The CM infinity exponents differ by one at every complex place,
so quadratic induction gives parallel weight two.  Thus the even
degree $[F:\mathbb Q]=8$ creates no obstruction to
\citet[Theorem~B(b), p.~593]{Nekovar2012}.

For $\delta=1$, \citet[Theorem~3.13, author pp.~17--18]{ShuExpanded2021} gives
\[
 w(A_i)=-\left(\frac{rst\,d}{17}\right).
\]
The archimedean factor is $(-1/17)=1$ and the other finite local
factors are one.  Since the three residues $rst\,d$ are respectively
$1,7,4$ modulo $17$, the signs are $-1,+1,-1$.  This also agrees
with the original table of \citet[p.~219]{GrossRohrlich1978}.  The functional
equation forces the first and third central values to vanish.

For the remaining factor, \citet[Theorem~2.1, p.~208]{GrossRohrlich1978}
gives a rational point of infinite order in $A_2(\mathbb Q)$.
If its central value were nonzero, the Hilbert form just described
would satisfy every hypothesis of Nekov\'a\v r Theorem~B(b).
That theorem would give zero rational Bloch--Kato Selmer group
for each coefficient realization.  But its representation induces
to $V_\ell(A_2)$ over $\mathbb Q$, by the Jacobi-sum construction;
Shapiro's lemma would therefore give
$H^1_f(\mathbb Q,V_\ell(A_2))=0$.
The Kummer class of the rational point of infinite order contradicts
this conclusion.  Hence the middle central value vanishes as well.
Its root number is positive, so its analytic order of vanishing is
at least two.
\end{proof}

\begin{corollary}[No rank-zero abelian quotient at seventeen]
\label{cor:seventeen-no-rankzero-quotient}
Every positive-dimensional abelian quotient over $\Q$ of
$\Jac(F_{17})$ has positive Mordell--Weil rank over $\Q$.
\end{corollary}
\begin{proof}
\citet[Concluding remark~2, p.~665]{McCallum1988} states
$\mathbb Z[\zeta_{17}]$-rank one over $K$ for the two wild-split
classes at exponent seventeen, namely the
classes of $(1,1,15)$ and $(1,3,13)$; the remaining class is the
wild-nonsplit class $(1,2,14)$. The residues $-2rst\,d$ are
$15,3,9$ modulo $17$, respectively, identifying those two split classes;
also $17\nmid B_8B_{10}$. This gives rational rank one:
the vector space $M=A_i(K)\otimes_{\mathbb Z}\mathbb Q$ is
one-dimensional over $K$, and $\operatorname{Gal}(K/\mathbb Q)$
acts semilinearly through its natural action on the CM
endomorphisms.  Galois descent (equivalently Hilbert~90 in dimension
one) gives $\dim_{\mathbb Q}M^{\operatorname{Gal}(K/\mathbb Q)}=1$;
these invariants equal $A_i(\mathbb Q)\otimes\mathbb Q$, since
taking invariants commutes with tensoring by $\mathbb Q$ for a
finite group.  Together with Gross--Rohrlich's
point of infinite order on $A_2$, this shows that every absolutely
simple factor in the displayed isogeny has positive rational rank.
Consequently no positive-dimensional abelian quotient of
$\operatorname{Jac}(F_{17})$ has rank zero over $\mathbb Q$.
Indeed a simple isogeny factor of such a quotient must occur in
the displayed decomposition, and rank is additive and isogeny
invariant.  This final rank statement uses McCallum's published
conclusion; the central-value obstruction above does not need it.
\end{proof}

\paragraph{Scope of the obstruction.}
The calculation verifies the finite-place hypothesis rather than
leaving it as an unresolved arithmetic condition.  The failed input
is precisely central nonvanishing: all three values are zero.
It rules out a rank-zero abelian target obtained from the Fermat
Jacobian and the present rank-zero Selmer theorem.  It proves no
noncuspidal section and no failure of a section conjecture.  A
rank-one or rank-two descent calculation together with a curve-level
local intersection argument would be a different route; none is
asserted here.

The section classification at exponent $17$ is nevertheless given
by Theorem~\ref{thm:regular-solvable-intro}; its proof uses the
regular-prime local argument rather than a rank-zero quotient.
